\section{Introduction}\label{sec:introduction}

High-dimensional data have emerged as a fundamental challenge in classification problems across diverse domains, including genomics, text mining, and computer vision. These datasets typically contain hundreds or even tens of thousands of features, while only a limited number of labeled samples are available, resulting in the well-known \(p \gg n\) scenario. Under such conditions, conventional learning algorithms often suffer from severe overfitting, increased computational complexity, and reduced interpretability due to the curse of dimensionality~\cite{lili2017,guyon2003}. Feature selection (FS) has therefore become an essential preprocessing technique that aims to identify a compact and informative subset of features while preserving classification performance. However, selecting an optimal feature subset remains highly challenging because the search space grows exponentially with the number of features \(D\), and the underlying subset-selection problem is NP-hard~\cite{liang2025}. Moreover, individual feature relevance measures are insufficient to characterize complex feature interactions and dependencies, which are often critical for achieving robust classification performance.

FS methods are generally categorized into filter, wrapper, and embedded approaches~\cite{kohavi1997,lili2017}. Filter methods rank features according to statistical criteria and offer high computational efficiency, but they usually evaluate features independently of the classifier and may overlook feature interactions relevant to classification. Embedded methods incorporate feature selection into the learning process, providing a compromise between selection efficiency and predictive performance; however, their effectiveness is often dependent on the characteristics of specific learning models. Wrapper methods directly optimize feature subsets according to classifier performance and can achieve competitive classification results, but the repeated model evaluations make them computationally expensive, especially in high-dimensional search spaces. To address the combinatorial complexity of wrapper-based FS, evolutionary computation has been extensively investigated as a promising optimization paradigm~\cite{liang2025,kennedy1995,shi1998,kennedy1997}. By maintaining a population of candidate solutions, evolutionary algorithms can perform global exploration of feature subsets without requiring gradient information. Among them, particle swarm optimization (PSO), genetic algorithms, and differential evolution have demonstrated promising results on high-dimensional FS tasks.

PSO has become one of the most widely adopted evolutionary algorithms for FS due to its simple structure, efficient information-sharing mechanism, and fast convergence. Nevertheless, conventional PSO-based FS methods still encounter three structural limitations when dealing with sparse and high-dimensional search spaces. First, random initialization distributes particles without prior guidance throughout the binary search space, whereas effective feature subsets usually occupy only a small region with limited active dimensions. Consequently, a considerable number of fitness evaluations may be wasted on redundant and suboptimal solutions before the swarm identifies promising search regions. Second, conventional uniform feature perturbation strategies treat all dimensions equally and neglect the feature-level improvement information accumulated during previous evaluations, preventing the search from exploiting historically informative features. Third, existing PSO-based FS methods typically employ fixed objective weights to balance classification error and feature subset size during the entire optimization process. However, the desired trade-off between feature reduction and accuracy refinement varies across search stages: aggressive sparsity enhancement is preferred in early exploration, whereas accurate decision boundary refinement becomes more important after a compact subset has been obtained.

Recent evolutionary FS studies have attempted to improve search performance through enhanced population initialization, adaptive evolutionary operators, and refined fitness evaluation strategies~\cite{wangjia2022,wang2022,xuezhang2024,wei2025}. Although these approaches have achieved promising results, they usually focus on isolated components of the optimization process and do not jointly address inefficient initialization, insufficient exploitation of accumulated feature knowledge, and inflexible objective weighting. Recently, multi-task PSO has emerged as an effective paradigm for high-dimensional FS with fixed-length binary representations. For instance, MEL~\cite{mel2024} adopts cooperative subpopulations and a fixed-length representation to efficiently handle feature spaces containing tens of thousands of dimensions. Despite its effectiveness, existing multi-task PSO-based FS frameworks still leave several challenges unresolved: the initial population is generally generated without sufficient feature-level prior knowledge, useful feature-importance information discovered during evolution is not explicitly retained and exploited, and the balance between classification performance and feature sparsity remains unchanged across different search stages.

To address these limitations, this paper proposes Knowledge-Guided Dual-Phase Swarm Evolutionary Learning (KDS-EL), a closed-loop framework that cohesively integrates domain prior injection, cross-iteration knowledge accumulation, and dynamic objective scheduling into a cooperative multi-task PSO backbone.

The main contributions of this paper are summarized as follows:
\begin{enumerate}
\item \textit{MI-Guided Initialization with Adaptive Re-initialization}: A hybrid seeding strategy that injects continuous feature-relevance priors into population construction to eliminate early-stage search waste, coupled with an adaptive re-initialization mechanism that restores swarm diversity upon evolutionary stagnation (Section~\ref{sec:mi}).
\item \textit{Knowledge-Guided Elite Refinement (KER)}: A targeted elite perturbation operator that dynamically harvests cross-iteration feature-importance feedback with exponential memory decay, directing elite solutions toward historically effective features within a strictly bounded evaluation overhead (Section~\ref{sec:ker}).
\item \textit{Two-Phase Weight Scheduling (TPS)}: A stage-aware objective scheduling mechanism that temporally decouples early aggressive feature pruning from late-stage accuracy refinement, overcoming the rigidity of static weighting (Section~\ref{sec:twophase}).
\item \textit{Comprehensive Empirical Validation}: Extensive evaluations across 20 high-dimensional benchmark datasets (spanning 500 to 54{,}675 features) against state-of-the-art single-objective competitors and precursor frameworks, rigorously validated through non-parametric statistical testing, component contribution analysis, and parametric sensitivity investigations (Section~\ref{sec:results}).
\end{enumerate}

Extensive experiments demonstrate that KDS-EL achieves the highest classification performance, obtaining an average classification accuracy of 0.9574 with only 87.8 selected features, corresponding to an average feature selection ratio of merely 1.1\%. Compared with recent top-tier competitor bIRIME~\cite{su2024} and established competitor FWPSO~\cite{wangjia2022}, KDS-EL improves classification accuracy by 1.20 and 3.97 percentage points, while reducing the selected feature count by 78.5\% and 91.8\%, respectively. Furthermore, KDS-EL achieves a 3.4$\times$ speedup over the foundational multi-task precursor MEL~\cite{mel2024}, reducing average runtime from 175.0\,s to 51.3\,s. Pairwise Wilcoxon signed-rank tests with Holm--Bonferroni step-down correction confirm that the performance advantages are statistically significant over all 11 comparison algorithms and ablative variants.

The remainder of this paper is organized as follows. Section~\ref{sec:related} formulates the feature selection optimization problem and reviews related literature and background concepts. Section~\ref{sec:proposed} presents the proposed KDS-EL framework and formulates its three constituent mechanisms in detail. Section~\ref{sec:results} presents comprehensive experimental results, statistical tests, component contribution analysis, and parameter sensitivity investigations. Finally, Section~\ref{sec:conclusion} concludes the paper and discusses prospective research directions.
